The complete menu experiment gives a positive economic role for the learned
continuation at a stated computational budget. All 640 method fits and
1,920 stage candidates are retained; no fit failed and no stage used the
fallback. Table~\ref{tab:r16menucomparisons} reports every final-stage direct
comparison, pooling the sixteen declared streams and all 384 task-state
queries within each of the eight cells. Its units are $10^{-4}$, the fixed
materiality and equivalence margin. The intervals belong to the same
216-event family as the complete earlier-stage record.

\begin{table}[!htbp]
\centering\footnotesize
\caption{Incremental Value of the Reusable NBO Continuation}\label{tab:r16menucomparisons}
\begin{tabular}{llrrrr}
\toprule
Economy & $d$ & NBO $-$ Vector & NBO $-$ Raw actor & NBO $-$ DPO & NBO $-$ SAA \\
\midrule
Original & 10 & [0.00, 0.01] & [0.06, 0.07] & [0.06, 0.07] & [-0.01, 0.01] \\
Original & 50 & [0.03, 0.05] & [0.25, 0.26] & [0.24, 0.25] & [-0.03, -0.02] \\
Quarterly & 10 & [-0.02, 0.02] & [0.20, 0.26] & [0.12, 0.17] & [-0.06, -0.01] \\
Quarterly & 50 & [0.22, 0.27] & [0.82, 0.90] & [0.76, 0.83] & [-0.41, -0.36] \\
Long & 10 & [0.14, 0.92] & [5.67, 6.70] & [-1.03, -0.19] & [-1.99, -1.22] \\
Long & 50 & [16.38, 18.13] & [18.02, 19.72] & [3.27, 4.25] & [-2.68, -1.89] \\
Intermediate & 10 & [-0.01, 0.03] & [0.25, 0.33] & [0.15, 0.21] & [-0.03, 0.02] \\
Intermediate & 50 & [0.66, 0.79] & [1.47, 1.63] & [1.30, 1.44] & [-0.49, -0.39] \\
\bottomrule
\end{tabular}
\begin{minipage}{0.97\linewidth}\footnotesize\emph{Notes:} Intervals are in units of $10^{-4}$ and use the final, third work stage. Every cell averages all sixteen declared streams and all 384 fixed current-period queries. Each query has 512 independent antithetic future pairs; opposite signs form one observation. Five methods share each pair. All three stages and all eight economic-dimension cells contribute to one 216-event family with error probability $0.018$. Current reward and a known terminal quadratic are removed before the random residual is bounded. The point estimator underlying each interval is the clipped residual estimator plus this known offset. Numerical arithmetic, clipping and the unclipped Gaussian law are included. The materiality margin is $10^{-4}$.
\end{minipage}
\end{table}


In both previously unpiloted fifty-dimensional designs, NBO yields materially
higher mean payoff than the materialized Raw and direct-policy actors. Its
lower gain over the Raw actor is $0.001802$ in Long and $0.000147$ in Intermediate;
the corresponding lower gains over direct policy are $0.000327$ and
$0.000130$, all rounded downward. Each exceeds $10^{-4}$. In Long at
dimension fifty, NBO also exceeds the fitted vector costate by at least
$0.001638$. These comparisons measure the economic consequences of using the
complete scalar continuation procedure. They do not estimate an independent
critic-risk difference or isolate variance removal from nonlinear
approximation and optimization.

\begin{table}[!htbp]
\centering\footnotesize
\caption{Complete Construction and Query Work}\label{tab:r16menuwork}
\begin{tabular}{llrrrrr}
\toprule
Economy & $d$ & NBO & Vector & Raw actor & DPO & Cached SAA \\
\midrule
Original & 10 & 12.52 & 8.73 & 6.86 & 16.21 & 15.60 \\
Original & 50 & 12.67 & 9.85 & 7.07 & 20.24 & 27.68 \\
Quarterly & 10 & 10.36 & 7.69 & 6.09 & 21.85 & 25.23 \\
Quarterly & 50 & 14.34 & 11.15 & 8.28 & 38.93 & 62.90 \\
Long & 10 & 12.05 & 8.87 & 6.96 & 26.53 & 30.42 \\
Long & 50 & 14.59 & 11.06 & 8.16 & 38.03 & 67.04 \\
Intermediate & 10 & 11.96 & 9.16 & 7.01 & 36.55 & 44.71 \\
Intermediate & 50 & 12.70 & 9.62 & 7.18 & 45.58 & 79.46 \\
\bottomrule
\end{tabular}
\begin{minipage}{0.97\linewidth}\footnotesize\emph{Notes:} Seconds per stream, averaged over all sixteen streams, including unsuccessful fits. Each entry is the actual outer-process clock for all three stages, imports, source checks, reusable training data, optimizer and derivative work, all saved queries and durable output. Fixed dependency installation and source extraction precede the clock. Independent confirmation is a shared measured bundle, reported separately in the complete account; the table does not allocate one-fifth of that bundle as an observed standalone verification cost.
\end{minipage}
\end{table}


The payoff improvements over direct policy in those two fifty-dimensional
cells accompany lower measured construction and query work. NBO's mean
complete three-stage clocks are $14.59$ seconds against $38.03$ for direct
policy in Long, and $12.70$ against $45.58$ in Intermediate. The materialized
Raw actor and vector predictor have lower construction and query clocks
than NBO in every final cell. Thus the economic increment over these two
representations has a computational cost. The comparison with direct policy
instead improves both the protected payoff and the observed work at the
stated final budget in the two specified cells.

Cached sample-average optimization supplies a further comparison with direct
simulation. In the six final cells outside Long, the complete NBO--SAA
interval lies within $[-10^{-4},10^{-4}]$. NBO uses less measured construction
and query work than the final sixty-four-path SAA in every one of these
cells. The reductions range from $19.73$ to $84.02$ percent of the SAA mean
clock. This is a comparison of the prespecified final-stage procedures. It
does not identify the minimum work required for a given accuracy: the
complete record also retains the cheaper four- and sixteen-path SAA stages.

The Long calibration displays the corresponding accuracy tradeoff. SAA has
a protected advantage exceeding $10^{-4}$ over NBO in both dimensions, at
higher final construction and query work. In Long at dimension ten, direct
policy also has a positive protected advantage over NBO, although its
material superiority at $10^{-4}$ is unresolved. Across all three stages,
NBO has nine material lower endpoints against Vector, seven against the Raw
actor, four against direct policy, and none against SAA, out of twenty-four
comparisons with each. The respective numbers of strictly negative NBO
upper endpoints are zero, three, seven, and nineteen. The complete intervals
in technical Section~\ref{app:r16menurecord} retain every reversal and every
unresolved result. Positive differences and practical equivalence are
reported as overlapping conclusions when both inequalities hold.

Independent confirmation is a shared computation. Its complete mean bill
ranges from $10.73$ to $82.20$ seconds per stream across the eight cells;
Table~\ref{tab:r16menusharedwork} gives every cell and the full clock range.
Adding that same bill to each method preserves the final-cell work
ordering, but the resulting total is an explicit conservative allocation of
the complete five-method, three-stage audit. It is not the time of an
independently run single-method verifier. The timing comparisons describe
the measured finite collection of runs and carry no unobserved-runtime
confidence claim.

\begin{table}[!htbp]
\centering\footnotesize
\caption{Conditional Accuracy on the Prescribed Scalar Intervention}\label{tab:r16menuscalar}
\begin{tabular}{llrrr}
\toprule
Economy & $d$ & Gap $\leq10^{-4}$ & Smallest gap bound & Largest gap bound \\
\midrule
Original & 10 & 16/16 & 0.0000002 & 0.0000002 \\
Original & 50 & 16/16 & 0.0000002 & 0.0000004 \\
Quarterly & 10 & 16/16 & 0.0000016 & 0.0000024 \\
Quarterly & 50 & 16/16 & 0.0000138 & 0.0000141 \\
Long & 10 & 0/16 & 0.0440599 & 0.0454588 \\
Long & 50 & 0/16 & 0.1962975 & 0.1963168 \\
Intermediate & 10 & 16/16 & 0.0000023 & 0.0000024 \\
Intermediate & 50 & 0/16 & 0.0002819 & 0.0002846 \\
\bottomrule
\end{tabular}
\begin{minipage}{0.97\linewidth}\footnotesize\emph{Notes:} All 128 conditional statements are retained under a separate error probability of $0.002$, with 65,536 independent antithetic pairs for each fixed scalar candidate. Each upper bound covers the whole continuous scalar action segment at the predeclared fixed high-dimensional state. It is not a bound on the complete vector action problem or continuous-time adapted policies. Endpoints are rounded upward.
\end{minipage}
\end{table}


The scalar exercise provides a separate absolute-accuracy result. Eighty of
the 128 fixed candidates have a protected gap below $10^{-4}$ over their
entire continuous scalar intervention interval. All sixteen streams attain
this criterion in Original and Quarterly at both dimensions and in
Intermediate at dimension ten. Long at both dimensions and Intermediate at
dimension fifty retain larger gap upper bounds. Nonattainment of an upper
certificate does not establish a corresponding lower bound on actual
suboptimality. These statements concern the fixed current scalar
intervention with the common future reference policy; they do not assert
near-optimality over all vector actions or adapted continuous-time controls.
